% !PRINTED 507
\weberpage{507}

In order to be able to regard $\varkappa(\omega)$ itself as a class invariant, one
applies the \gesperrt{Gauss} transformation (\S\,9, \S\,32). According to
it:

\[ \varkappa^2\left(\dfrac{\omega}{2}\right) = \dfrac{4\varkappa(\omega)}{[1+\varkappa(\omega)]^2}, \]

from which:

\[ \varkappa(\omega) = \dfrac{\varkappa^2\left(\dfrac{\omega}{2}\right)[1+\varkappa^2(\omega)]}{4-2\varkappa^2\left(\dfrac{\omega}{2}\right)}. \]

Thus $\varkappa(\omega)$ is expressed rationally by $\varkappa^2(\omega)$ and
$\varkappa^2\left(\dfrac{\omega}{2}\right)$, and $\omega' = \dfrac{\omega}{2}$ satisfies the equation

\weq{8}{4A\omega'^2 + 2B\omega' + C = 0.}

If $C$ is an even number, which we may assume when
$D\equiv 0\ (\mathrm{mod}\ 4)$ or $\equiv 1\ (\mathrm{mod}\ 8)$, then the
factor 2 can be cancelled in (8), and the discriminant of (8) is
equal to $D$. But if $C$ is odd, which is necessary when $D\equiv 5\ (\mathrm{mod}\ 8)$, then the discriminant of (8) is equal to $4D$. From this
it follows by 2. and 3.:

4. a) If $D\equiv 0\ (\mathrm{mod}\ 4)$, then $\varkappa(\omega)$ is \gesperrt{rationally expressible by the class invariants of the discriminant $4D$.}

b) If $D\equiv 1\ (\mathrm{mod}\ 8)$, then $\varkappa(\omega)$ is \gesperrt{expressible by $\sqrt{D}$ and by the class invariants of the discriminant $4D$.}

c) If $D\equiv 5\ (\mathrm{mod}\ 8)$, then $\varkappa(\omega)$ is \gesperrt{rationally expressible by $\sqrt{D}$ and by the class invariants of the discriminant $16D$.}

\begin{center}
\S\,136.\quad \gesperrt{Quadratic Degrees of Transformation.}
\end{center}

When the degree of transformation $n$ is an \gesperrt{odd square}
number, then by \S\,72, 5. one has for a variable $\omega$:

\weq{1}{M = (-1)^{\frac{a-1}{2}} \partial^3\left(\dfrac{\eta\left(\dfrac{c+\partial\omega}{a}\right)}{\eta(\omega)}\right)^6,}
%CONT>>
